\section{Training setup}
\label{sec:training}

\subsection{Data splits and normalization}
\label{subsec:train-splits}
All models are trained on the processed datasets described in Section~\ref{sec:data-generation}. As explained there, the training, validation, and test scenarios come from separately generated pools with separate fixed dataset seeds rather than from a random partition of the pool, which guarantees that no scenario is shared between them. The main experiment uses $10{,}000$ training scenarios, $1{,}000$ validation scenarios, and $2{,}500$ test scenarios per solver target, with the same scenario design reused across the hydrostatic, MUSCL-HR, and Boussinesq labels so the targets are directly comparable.

Normalization statistics are estimated on the training split only and then applied unchanged to the validation and test splits, so that no test information leaks into the feature scaling. Each active input channel and the trajectory target are standardized to zero mean and unit variance; the physical-unit metrics reported in Section~\ref{subsec:eval-metrics} are obtained by reversing this target standardization at evaluation time.

\subsection{Optimization}
\label{subsec:train-optim}
All models are trained with the AdamW optimizer at an initial learning rate of $5 \times 10 ^ {-4}$ and a weight decay of $10 ^ {-6}$. The learning rate follows a cosine-annealing schedule that decays to a floor of $5 \times 10 ^ {-5}$ over the training budget, and the gradients are clipped to a maximum norm of $1.0$ for stability. We use a batch size of $64$.

Because the target is a $50$-frame trajectory, we train with a horizon-weighted mean-squared error that increases the weight of later frames,
\[
\mathcal{L} = \frac{\sum_{t = 1} ^ {T} w_t\, \mathrm{MSE}_t}{\sum_{t = 1} ^ {T} w_t},
\qquad
w_t = w_{\min} + (w_{\max} - w_{\min})\left(\tfrac{t - 1}{T - 1}\right) ^ {p},
\]
where $\mathrm{MSE}_t$ is the mean-squared error of frame $t$ and we use $w_{\min} = 0.8$, $w_{\max} = 2.0$, and $p = 1.7$. Without this weighting, the optimizer tend to fit the easy early frames and leave the later, lower-amplitude part of the rollout blurry, which is exactly the regime that matter for a forecast horizon. The relative $L_2$ error on the validation split is used both as the checkpoint-selection metric, keeping the best epoch, and as the early-stopping criterion with a patience of $20$ epochs. One detail worth stating: this selection metric is the relative $L_2$ computed on the \emph{normalized} validation targets, whereas the accuracy tables we report later are in physical units after the normalization is reversed. Because standardization subtracts the channel mean, the normalized and physical relative $L_2$ have different denominators and are not the same number, so the best epoch is chosen on the normalized criterion and then scored on the physical metrics rather than being selected directly on the physical error. To keep the comparison fair, all four models share a single training budget: every model is allowed up to $256$ epochs, the largest budget used by any model, and the same optimizer, schedule, loss, and stopping rule applied to all of them unless a model-specific configuration explicitly changes only the architecture. Early stopping then halts each model on its own, when its validation error stops improving, so a faster-converging baseline is not penalized and a slower model is not cut short. In practice the budget is an upper bound rather than a fixed number of epochs, and most models stop well before reaching it.

\subsection{Experiment configurations}
\label{subsec:train-configs}
Each model--target pair is defined by a single configuration file. The primary configurations train the Fourier Neural Operator on the hydrostatic, MUSCL-HR, and Boussinesq labels, and the hydrostatic comparison set additionally includes FFNO, FNO-modes8, FNO-modes20, U-FNO, WNO, CNN, and U-Net models under the matched protocol above. Hydrostatic is the main reference target, while MUSCL-HR and Boussinesq are the main comparison targets; Boussinesq is included because it gives a weakly dispersive baseline on the same scenarios, not because it is treated as final physical truth.

The windowed configurations use the same hydrostatic data but change the supervised task from a direct $50$-frame forecast to a $K = 5$ frame block forecast conditioned on recent surface-elevation frames. The reported Window-FNO-5 and Window-FFNO-5 rollouts are seeded with the true first frame at evaluation time, so they are used as drift diagnostics and are not treated as identical to the direct single-pass prediction problem.

For the uncertainty experiments we train a deep ensemble of seven members that differ only in their random seed, following the procedure of Section~\ref{subsec:uncertainty}. The ensemble shares the architecture and the training protocol of a single FNO model and only changes the seed, so any spread between member reflects optimization stochastically rather than a different model design.

\subsection{Hardware and runtime setting}
\label{subsec:train-hardware}

Dataset generation was performed on an 8-core AMD EPYC 9B45 CPU server, because the numerical reference solvers are the main computational bottleneck during benchmark construction. In particular, each scenario must be advanced by explicit numerical time integration before it can become a supervised training label. Model training and evaluation were then run on accessible local hardware: an Intel Core i5-13420H CPU and an NVIDIA GeForce RTX 4050 Laptop GPU. We report this setting explicitly because one goal of the benchmark is to make neural-surrogate experimentation feasible for smaller laboratories, student researchers, and local communities without access to large GPU clusters.

The resulting runtime numbers should therefore be interpreted as measurements under a resource-conscious research setup rather than as optimized production benchmarks. They quantify the acceleration of the neural surrogate relative to the reproducible reference implementations used in this study, not relative to highly optimized operational tsunami models or GPU-accelerated numerical solvers. A global random seed of $42$ is set for the Python, NumPy, and PyTorch random number generators, the ensemble members use the fixed seed $\{11, 22, 33, 44, 55, 66, 77\}$, and the data splits are fixed by construction as described above. Every run writes its resolved configuration, its full epoch history, and its best and last resolved checkpoints into a per-run directory. The speed-evaluation outputs also record the selected device, Python/PyTorch/CUDA versions, precision label, and TF32 backend state, so each reported runtime member can be traced back to the configuration and hardware that produced it.
